\section*{Introduction to factorial}
The factorial of natural number $r$ is denoted by $r!$

\subsection*{Definition}
Factorial of natural number $r$ is the product of all natural number from 1 to r.
\begin{enumerate}[label=(\alph*)]
    \item $5! = 1\cdot2\cdot3\cdot4\cdot5 = 120$
    \item $3! = 1\cdot2\cdot3 = 6$
\end{enumerate}

\subsection*{Observation}
\begin{enumerate}
    \item Factorial of fraction or negative numbers are not defined.
    \begin{enumerate}
        \item $(\frac{3}{4})! = undefined$
        \item $(-12)! = undefined$
    \end{enumerate}
    \item Factorial of every natural number is $even$ except for 1.
\end{enumerate}
\subsection*{Common Factorial Values}

\begin{table}[h!]
\centering
\begin{tabular}{c c}
\hline
$n$ & $n!$ \\
\hline
1 & 1 \\
2 & 2 \\
3 & 6 \\
4 & 24 \\
5 & 120 \\
6 & 720 \\
\hline
\end{tabular}
\caption{Common factorial values}
\label{table:factorial}
\end{table}

\subsection*{Some Facts}
\begin{enumerate}
    \item Factorial of number greater than or equal to 5 will always have $0$ in their unit place.
    \item If a number is perfect square then only 6 of these are possible in unit place \{0,1,4,5,6,9\}
    \item Sum of factorials after 4 will have unit place digit as 3.
\end{enumerate}

\subsection*{Question: What is the minimum value of r if $r!$ is divisible by 13 where $r\in N$}
\textbf{Solution:} $r! = 1\cdot2\cdot3\cdots r$\\
In order to a number to be divisible by 13 it should have 13 in is factors, the only point the 13 is included in factorial of number when $r = 13$. So, $r_{min} = 13$.

\subsection*{Question: Is $13!$ divisible by 26}
\textbf{Solution: } Both 13 and 2 comes in $13!$. Also the factor of 26 is 2 and 13 hence, 13 is completely divisible by 26.

\subsection*{Question: Find the minimum value of $r(r \in N)$ if $r!$ is divisible by 16.}
\textbf{solution: } \\
$1 \times2 \times3 \times4 \times5 \times6 \times7 \times8 \cdots$ \\
Till here we have four 2's. 
From 2(1), 4(2),6(1), total four 2's we have.\\
And $16=2\times2\times2\times2$
Hence $r_{min} = 6$ which is completely divisible by 16.

\subsection*{Question: Solve for $x,y$ where $x,y \in N$ \\ $1!+2!+3!+\cdots+x! = y^2$ }
\textbf{Solution: } 
$r!$ will have unit digit as $0$ for $r \geq 5$. Also the sum of factorial natural number for $r \geq 4$ will have unit digit 3. We know any number which is perfect square will not have unit place digit 3. Hence the equation will only satisfies for $x=1$ and $x = 3$.

\subsection*{Question: Find the $x(\in N)$\\ such that it satisfy. $1!+2!+3!+\cdots+x!= 9\cdot y$}
\textbf{Solution: }
The equation to satisfy the LHS should be multiple of 9.\\
Now the sum of factorial till $5$ is divisible by 9. \\
$1!+2!+3!+4!+5! = 153$\\
The every factorial of number $\geq 6$ is always be divisible by 9. As 3 comes two times till $6!$. If a number is divisible by 9 and we add a number which is also divisible by 9 then the resulting number will always be divisible by 9.
\newpage
\subsection*{Question: Evaluate the statement given below}
Let \[
x = \sum_{r=1}^{2n}r!
\]
\[
y = \sum_{r=1}^{n}(2r)! 
\]
Statements
\begin{enumerate}[label = (\alph*)]
    \item x is divisible by y
    \item y is divisible by x
\end{enumerate}
\textbf{Solution: }\\
$x = 1!+2!+3!+\cdots+(2n)!$ \\
$y = 2!+4!+6!+\cdots+(2n)!$ \\
Number of term in $x = 2n$ and\\
Number of term in $y = n$\\
Since number of term in y is less so y is never be divisible by x, that makes statement 1 false.\\
Now the x is odd, as after 2 all are even and when added 1 it become odd. And y is even because addition of factorial after 2 is even.\\
Odd can not divide even.
\newpage
\section*{Exponent}
\subsection*{Find the maximum value of $r(\in N)$ that $\frac{10!}{2^r}$ is an integer.}
\textbf{Solution: } \\
$10! = 1 \times2 \times3 \times4 \times5 \times6 \times7 \times8 \times9 \times10$\\
Number of two's: 
$(2, 1), (4,2), (6, 1), (8, 3), (10, 1  )$
So total 8 two's are there we can divide the $10!$. Hence $r_{max} = 8$

\section*{General Formula}
Exponent of prime $p$ in $k!$ is given as.
$\left[\frac{k}{p}\right]+\left[\frac{k}{p^2}\right]+\left[\frac{k}{p^3}\right]+\cdots$ \\
Where [.] denotes G.I.F

\subsection*{Question: Exponent of 2 in 10!}
\textbf{Solution: }\\
$\left[\frac{10}{2}\right]+\left[\frac{10}{4}\right]+\left[\frac{10}{16}\right]+\cdots$ \\ \\
$5+2+1 = 8$

\subsection*{Explanation}
Let say we are find the maximum power of 3 which can divide the 100!. \\
$1 \times2 \times3 \times4 \times5 \times6 \times7 \times8 \times9 \times \cdots \times27 \times \cdots \times 81\times \cdots \times99 \times100 $ \\
When we are doing $\left[\frac{100}{3}\right]$ we get 33 which convey a meaning that there are 33 number which have 3 multiple. This is not sufficient, because 81 is although divisible by 3 but it counted just one power , we will use 9  and then sub sequence power of 3 to calculate the maximum power which can divide the given number. Just ask one simple question how many power does a number contribute in 3

\subsection*{Question: Find $r_{max}$ if $120!$ is  divisible by $6^r$}
\textbf{Solution: }
\\
$6^r = 2^r\cdot 3^3$\\
We will find the $120!$ divisibility by $2^r$ and $3^r$. The common power or the lower one will do the work.

\subsection*{Question: Find $r_{max}$ for $12^r$ divides 120!}
Answer is 58

\subsection*{Question: Find the number of zeros at the end of $150!$}

\textbf{Solution:}

Since $10 = 2 \times 5$, we find the powers of $2$ and $5$ in $150!$.

\medskip
\textbf{Power of 2:}
\[
\begin{aligned}
\left[\frac{150}{2}\right]
+ \left[\frac{150}{4}\right]
+ \left[\frac{150}{8}\right]
+ \left[\frac{150}{16}\right]
+ \left[\frac{150}{32}\right]
+ \left[\frac{150}{64}\right]
+ \left[\frac{150}{128}\right]
&= 75 + 37 + 18 + 9 + 4 + 2 + 1 \\
&= 146.
\end{aligned}
\]

\textbf{Power of 5:}
\[
\begin{aligned}
\left[\frac{150}{5}\right]
+ \left[\frac{150}{25}\right]
+ \left[\frac{150}{125}\right]
&= 30 + 6 + 1 \\
&= 37.
\end{aligned}
\]

To form a factor of $10$, we need one $2$ and one $5$.  
Hence, the number of trailing zeros is

\[
\min(146, 37) = 37.
\]

\[
\boxed{37}
\]

\subsection*{Question: Let $57\cdot58\cdot59\cdots200=5^r\cdot I$ Where $r,T \in N$, the find the $r_{max}$.}
\textbf{Solution: } \\
\[
\begin{aligned}
57\cdot58\cdot59\cdots200 = \frac{1\cdot2\cdot3\cdots200}{1\cdot2\cdot3\cdots56}
&=\frac{120!}{56!}
\end{aligned}
\]
Here we have to find the power of 5. The $r_{max}=37$